% !TEX root = ../main.tex
%-----------------------------------------------------------------------
%  What has to be worked out, and how this direction relates to the other.
%  \input by ../main.tex -- not compiled on its own.
%-----------------------------------------------------------------------
\hd{What has to be worked out.}
Unordered, and certainly incomplete.

\smallskip
\begin{tabularx}{\textwidth}{@{}Y{0.5}Y{1.5}@{}}
\toprule
\textbf{Item} & \textbf{What is actually in the way} \\
\midrule
Input binaries, and \code{tile001.mitgrid} above all
  & \code{input\_binaries/} is untracked and produced outside this repository;
    nothing here regenerates it, and unlike SOMA there is no \code{gendata.py}
    to fall back on. The horizontal grid is not in the namelist at all ---
    \code{delX} and \code{delY} are commented out and
    \code{usingCurvilinearGrid = .TRUE.} reads \code{tile001.mitgrid} instead
    --- so that one 1.3\,MB file \emph{is} the grid. Producing a 1/4\dg\
    version of it is the first practical blocker, in concrete form. \\
\addlinespace
Grid generation
  & Whether the curvilinear grid is regenerated or interpolated, and whether
    the vertical stays at 36 levels. \\
\addlinespace
Spin-up strategy
  & Per section~2: the dominant cost, and the first decision. \\
\addlinespace
Decomposition and tile size
  & And whether the run fits the queue limits \code{PORTING.md} already flags
    for the 1\dg\ spin-up. \\
\addlinespace
Adjoint stability
  & Over the intended window, and what viscosity treatment it needs. \\
\addlinespace
Storage
  & The 1\dg\ five-year adjoint writes 8.3\,GB, and the spin-up holds 47\,GB of
    retained restarts inside an 85\,GB run directory. Both scale with the cell
    count. \\
\addlinespace
Cost-section indices
  & \code{code\_tap/cost\_atlantic\_heat.F} carries
    \code{isecbeg}/\code{isecend}/\code{jsec}/\code{kmaxdepth} as compiled-in
    \code{parameter} statements, so a new grid means recomputing them and
    rebuilding. A second argument for a change the surrogate direction already
    asks for: move them into \code{data.cost}. \\
\addlinespace
A machine
  & Whether sverdrup can carry this at all, or whether it forces the Perlmutter
    port that \code{PORTING.md} describes and nobody has validated. \\
\bottomrule
\end{tabularx}

\hd{Relation to the surrogate direction.}
That direction is scoped entirely at 1\dg, and nothing here changes it. The
connection runs one way: a surrogate that generalises across parameters at
1\dg\ is interesting; one that transfers to 1/4\dg\ would be the result worth
having. Its own limitations list already names the idealised configuration as a
thing not to assume away, though it says nothing about resolution --- which is
this direction's question rather than its. Whether the 1/4\dg\ configuration
ever becomes a training target or stays a validation target is not decided, and
does not need to be until the questions in section~1 have answers.

\hd{What this note is not.}
It is not a plan, not a resource request, and not a design. It is the arithmetic
that says whether writing those is worth anyone's afternoon, plus a list of what
would have to be settled first. The next thing to write is the spin-up decision
in section~2; everything else waits on it.
